\chapter{Complex Engineering Problems and Activities}
\label{ch:cep}

This chapter maps LectureAssist onto the Complex Engineering Problem (CEP) and
Complex Engineering Activity (CEA) attributes, showing where the project required
depth of knowledge, resolution of conflicting requirements, and analysis beyond the
routine application of a known method. Each attribute is addressed with reference to
the specific decision, defect or measurement in this report that exercises it, rather
than by general assertion.

\section{Complex Engineering Problems (CEP)}
\label{sec:cep}

{\footnotesize\singlespacing
\begin{longtable}{|c|p{3.4cm}|p{8.4cm}|}
\caption{Complex Engineering Problem attributes addressed by LectureAssist.}
\label{tab:cep}\\
\hline
\textbf{Attr.} & \textbf{Description} & \textbf{How the project addresses it} \\
\hline
\endfirsthead
\hline
\textbf{Attr.} & \textbf{Description} & \textbf{How the project addresses it} \\
\hline
\endhead
\hline
\endfoot
\hline
\endlastfoot

P1 & Depth of knowledge required (K3--K8) &
The system spans computer vision and image preprocessing for scene-adaptive OCR
(K3--K4), automatic speech recognition, transformer language models, embeddings and
vector retrieval (K4--K5), multi-agent software architecture and API/system design
(K5--K6), statistical evaluation design, LLM-as-judge, $F_1$ agreement between the
internal validator and an independent judge, calibration against a known-answer set
(K5, K8), and the research literature on question generation, chain-of-thought and
program-of-thought prompting, and agentic architectures (K8). \\
\hline
P2 & Range of conflicting requirements &
Several requirements pull directly against each other. \emph{Question quality vs.\
cost}: the agentic loop produces better-grounded, less repetitive and more
answer-correct questions but consumes roughly an order of magnitude more LLM calls
than a single-shot request. \emph{Model size vs.\ VRAM}: a larger generator would
produce better questions, but the generator, judge, Whisper, OCR and embedding models
must all coexist on one $16$\, GB T4. \emph{Yield vs.\ scientific honesty}: allowing
the $4$B generator to escalate to the $12$B model on malformed output would have
raised yield, but would have falsified the claim that the questions were written by
the $4$B, so escalation was disabled at a known cost in yield. \emph{Accept rate
vs.\ diversity}: a pipeline can maximise its accept rate by repeatedly producing one
good question, which forced the adoption of a duplicate-penalised effective
acceptance metric (Eq.~\ref{eq:effaccept}) in place of raw acceptance. \\
\hline
P3 & Depth of analysis required &
There is no unique correct design, and no off-the-shelf metric for ``is this a good
generated question''. Selecting a solution required analysing four generation
strategies (agentic, plain, CoT, PoT) against each other under a metric suite that
had itself to be designed, and it required identifying and eliminating confounds,
unequal content coverage between pipelines, a truncated summary that made whole
sections of the chapter unreachable by the planner, and silent model escalation,
any one of which would have invalidated the comparison. The central analytical
finding is that the four judged quality dimensions \emph{saturate} and therefore
discriminate nothing, so that the entire result rests on corpus-level automated
measures that a conventional quality-only evaluation would never have computed. That
conclusion is invisible without the analysis and inverts the naive reading of the
data. \\
\hline
P4 & Familiarity of issues &
The project integrates infrequently encountered components: EasyOCR with
scene-dependent preprocessing, Faster-Whisper, Gemma~3 served through Ollama,
Sentence-BERT embeddings, a FAISS index, an LLM-as-judge evaluation harness, SymPy
for machine-verified answer keys, and a Flask/ngrok deployment onto ephemeral
free-tier GPU infrastructure. \\
\hline
P5 & Extent of applicable codes &
No standard or code of practice exists for automatic educational question generation
or for the evaluation of LLM-generated assessment items. The evaluation protocol,
what to measure, how to verify an answer, how to calibrate the judge, had to be
designed from the research literature rather than taken from a standard. Where
published protocols did exist they had to be reconstructed from the papers
themselves: EQGBench has no public code release, and ReQUESTA publishes no judge
prompt at all because its rubric was applied by human raters. \\
\hline
P6 & Extent of stakeholder involvement &
Stakeholders include students (the primary users), instructors (whose lectures are
the input and whose assessments the tool must not undermine), institutions (which
hold copyright in lecture material and are accountable for academic integrity), and
the open-model community whose licences govern the deployed models. Their
requirements conflict: a student wants unlimited questions immediately, while an
instructor requires that the tool not present unverified answers as authoritative,
and an institution requires that copyrighted lecture material not be transmitted to a
third-party processor. \\
\hline
P7 & Interdependence &
The subsystems are tightly coupled: OCR and ASR quality bound the summary; the
summary bounds the planner's topic choice; the planner bounds question diversity; the
retrieval index bounds grounding; and the judge and Answer Generator bound the
credibility of every reported result. Two of the defects found during this project, a truncated summary and a capped content window, were failures of exactly this
interdependence: a bug several stages upstream degraded question diversity many
stages downstream, with no local symptom at the point of failure. The pipeline
labelling error disclosed in Section~\ref{sec:res-native} is a third instance, in
which a fault in the recording of provenance propagated into the interpretation of
every downstream measurement. \\
\hline
\end{longtable}
}

Table~\ref{tab:cep} shows that the project engages all seven CEP attributes, with P2
(conflicting requirements) and P3 (depth of analysis) being the most heavily
exercised. It is worth noting that P2 and P3 are linked in this project rather than
independent: almost every conflicting requirement listed under P2 was resolved by
\emph{measuring} the trade-off rather than by asserting a preference, and the
instruments built to perform those measurements are what generated the depth of
analysis recorded under P3.

\section{Complex Engineering Activities (CEA)}
\label{sec:cea}

{\footnotesize\singlespacing
\begin{longtable}{|c|p{3.4cm}|p{8.4cm}|}
\caption{Complex Engineering Activity attributes addressed by LectureAssist.}
\label{tab:cea}\\
\hline
\textbf{Attr.} & \textbf{Description} & \textbf{How the project addresses it} \\
\hline
\endfirsthead
\hline
\textbf{Attr.} & \textbf{Description} & \textbf{How the project addresses it} \\
\hline
\endhead
\hline
\endfoot
\hline
\endlastfoot

A1 & Range of resources &
The project draws on human resources (a three-member team with a supervisor), modern
computational tools (a GPU runtime, six distinct deep-learning models, a vector
index, a symbolic algebra system), open-licensed source material, and cloud compute, managed under a hard constraint of zero paid API spend on the generation path,
as documented in the budget of Section~\ref{sec:plan-budget}. \\
\hline
A2 & Level of interactions &
The work required interaction between the extraction, retrieval, generation,
evaluation and interface subsystems, each with its own failure modes, as well as
between team members under the responsibility allocation of
Table~\ref{tab:raci}, the project supervisor (whose requirements shaped the
evaluation design, notably the demands for CoT and PoT baselines and for treating
question generation and answer generation as separate cross-checked agents), and the
end users whose trust in the output determines whether the system is safe to use. \\
\hline
A3 & Innovation &
The project introduces three novel elements: scene-adaptive OCR that classifies each
frame before choosing its preprocessing, so that a lecture's \emph{board}: not
merely its audio, becomes usable content; a validator agent that selects
\emph{how} to remediate a rejected question (regenerate, fetch more context, or
replan the topic) rather than merely rejecting it; and an evaluation protocol that
separates question quality from answer correctness by having a distinct Answer
Generator agent re-solve every question blind, with the solver itself calibrated on a
machine-verified answer key. \\
\hline
A4 & Consequences to society and the environment &
The system widens access to effective retrieval practice for students who cannot pay
for cloud AI services, supporting UN SDG~4 (Quality Education) and SDG~10 (Reduced
Inequalities). Running small models locally is substantially less energy-intensive
per request than hyperscale cloud inference, and no model was trained. The principal
risk, a fluent question with a wrong marked answer, is analysed and bounded
rather than dismissed, and the report states explicitly that generated questions are
suitable for practice and not for graded assessment without human review.
Chapter~\ref{ch:impacts} develops these consequences in full. \\
\hline
A5 & Familiarity &
The project requires familiarity with computer vision, speech recognition, large
language models and prompting strategies, retrieval-augmented generation, agentic
architectures, experimental design and statistical evaluation, web API development,
and deployment onto constrained, ephemeral GPU infrastructure. \\
\hline
\end{longtable}
}

Table~\ref{tab:cea} demonstrates that the project satisfies all five Complex
Engineering Activity attributes.

\section{Mapping to Knowledge Profiles}
\label{sec:cep-knowledge}

Table~\ref{tab:knowledge} records which Washington Accord knowledge profiles the work
draws on and where in this report the evidence for each can be found, so that the
mapping can be checked against the substance of the project rather than accepted on
assertion.

{\footnotesize\singlespacing
\begin{longtable}{|c|p{5.0cm}|p{6.4cm}|}
\caption{Knowledge profile coverage and where each is evidenced in this report.}
\label{tab:knowledge}\\
\hline
\textbf{Profile} & \textbf{Description} & \textbf{Evidence in this report} \\
\hline
\endfirsthead
\hline
\textbf{Profile} & \textbf{Description} & \textbf{Evidence in this report} \\
\hline
\endhead
\hline
\endfoot
\hline
\endlastfoot

K3 & Engineering fundamentals &
Image preprocessing, contrast-limited adaptive histogram equalisation,
thresholding, polarity inversion, selected per frame class
(Section~\ref{sec:meth-ocr}) \\
\hline
K4 & Engineering specialist knowledge &
Transformer language models, speech recognition, sentence embeddings and
approximate nearest-neighbour search (Chapter~\ref{ch:background}) \\
\hline
K5 & Engineering design &
The five-agent Plan--Retrieve--Generate--Validate--Refine architecture and its
alternatives (Section~\ref{sec:meth-agentic}, Figure~\ref{fig:agentic-loop}) \\
\hline
K6 & Engineering practice &
Component selection against stated alternatives and rationale
(Table~\ref{tab:tools}); deployment on constrained infrastructure \\
\hline
K7 & Comprehension of the role of engineering in society &
Access, safety, privacy and environmental analysis (Chapter~\ref{ch:impacts}) \\
\hline
K8 & Research literature &
Two published protocols reconstructed from their source papers and applied with
declared deviations (Section~\ref{sec:res-protocol}) \\
\hline
\end{longtable}
}
